\section{Exact cofactor exposure on each physical row}
\label{sec:tpc61-cofactor-exposure}

We first isolate the part of the missing--high problem that is already
diagonal in the inherited row coefficients.  This must be done before any
residual or native grouping: after grouping, a statement about omitted
atomic mass is no longer monotone.  The notation below retains the literal
coefficient and raw counting normalization of TPC--53--TPC--60
\citep{WangTPC53,WangTPC54,WangTPC59,WangTPC60}.

\subsection{The realized cofactor measure}

Work in the identical-support specialization of
\cref{def:tpc61-identical-support}.  Thus a high Walsh atom is represented
precisely when its recovered pair \((m,r)\) belongs to the fixed retained
relation \(\cI_X\).  Let \(\cM_X\) be the declared source-row set and
retain the notation
\begin{equation}
 \cI_m=\{r\in\N:(m,r)\in\cI_X\},
 \qquad m\in\cM_X,
 \label{eq:tpc61-row-retained-cofactors}
\end{equation}
from \eqref{eq:tpc61-retained-cofactor-row}.  The set \(\cI_m\) is fixed
independently of the coefficient vector.

Every \(w\in\cW_H(y)\) has the unique recovered labels in
\eqref{eq:tpc61-recovered-mjp}--\eqref{eq:tpc61-recovered-rs}.  In
particular,
\begin{equation}
 n(w)=m(w)j(w)+h_0,
 \qquad
 p(w)=\Pplus(n(w))>y,
 \qquad
 r(w)=\frac{n(w)}{p(w)}.
 \label{eq:tpc61-recovered-prime-cofactor}
\end{equation}
Literal factorization is
\(b_w(\zeta)=\beta_w\zeta_{m(w)}\), as in
\eqref{eq:tpc61-literal-row-amplitude}; \(\beta_w\) contains every fixed
physical factor but not the bounded row coefficient \(\zeta_m\).

\begin{definition}[Realized row cofactor measure]
\label{def:tpc61-realized-cofactor-measure}
For \(m\in\cM_X\), define the finite positive measure on \(\N\)
\begin{equation}
 \nu_m(r)
 :=\sum_{\substack{w\in\cW_H(y)\\m(w)=m,\ r(w)=r}}
       |\beta_w|^2.
 \label{eq:tpc61-nu-m-r}
\end{equation}
Write
\begin{equation}
 W_H(m):=\sum_r\nu_m(r),\qquad
 W_R(m):=\sum_{r\in\cI_m}\nu_m(r),\qquad
 W_M(m):=\sum_{r\notin\cI_m}\nu_m(r).
 \label{eq:tpc61-row-HRM-weights}
\end{equation}
Thus \(W_H=W_R+W_M\) row by row.
\end{definition}

The measure \(\nu_m\) is not counting measure on possible cofactors.  It
counts only cofactors actually realized by the largest-prime factorization
of \(mj+h_0\), with the complete coefficient-free physical square attached.

\begin{proposition}[Exact row disintegration]
\label{prop:tpc61-exact-row-disintegration}
For every row-coefficient vector \(\zeta\),
\begin{align}
 \cD_H(\zeta)
 &=\sum_mW_H(m)|\zeta_m|^2,
 \label{eq:tpc61-DH-row-disintegration}\\
 \cD_R(\zeta)
 &=\sum_mW_R(m)|\zeta_m|^2,
 \label{eq:tpc61-DR-row-disintegration}\\
 \cD_M(\zeta)
 &=\sum_mW_M(m)|\zeta_m|^2.
 \label{eq:tpc61-DM-row-disintegration}
\end{align}
In particular, the represented/missing split is an orthogonal coordinate
projection before residual or native grouping.
\end{proposition}

\begin{proof}
The sets \(\cW_R\) and \(\cW_M\) are the disjoint subsets of
\(\cW_H\) on which respectively \(r(w)\in\cI_{m(w)}\) and
\(r(w)\notin\cI_{m(w)}\).  Substitute
\eqref{eq:tpc61-literal-row-amplitude} into the three raw atomic
squares and group nonnegative terms first by \(m\), then by \(r\).
No coherent sum occurs in this calculation.
\end{proof}

\subsection{Sharp uniform and distinguished exposure constants}

For every row with \(W_H(m)>0\), define
\begin{equation}
 \varepsilon_m
 :=\frac{W_M(m)}{W_H(m)},
 \qquad
 \rho_m:=\frac{W_R(m)}{W_H(m)}=1-\varepsilon_m.
 \label{eq:tpc61-row-defect-and-exposure}
\end{equation}

\begin{theorem}[Sharp rowwise exposure minimax]
\label{thm:tpc61-rowwise-exposure-minimax}
Assume the bounded coefficient class contains a nonzero multiple of every
coordinate direction.  If \(\cD_H\) is not identically zero, then
\begin{equation}
 \boxed{
 \sup_{\cD_H(\zeta)>0}
 \frac{\cD_M(\zeta)}{\cD_H(\zeta)}
 =\max_{m:W_H(m)>0}\varepsilon_m,}
 \label{eq:tpc61-sharp-missing-minimax}
\end{equation}
and
\begin{equation}
 \boxed{
 \inf_{\cD_H(\zeta)>0}
 \frac{\cD_R(\zeta)}{\cD_H(\zeta)}
 =\min_{m:W_H(m)>0}\rho_m
 =1-\max_{m:W_H(m)>0}\varepsilon_m.}
 \label{eq:tpc61-sharp-represented-minimax}
\end{equation}
Both identities are sharp at every finite \(X\).
\end{theorem}

\begin{proof}
Whenever \(\cD_H(\zeta)>0\),
\begin{equation}
 \frac{\cD_M(\zeta)}{\cD_H(\zeta)}
 =\sum_{m:W_H(m)>0}
   \frac{W_H(m)|\zeta_m|^2}{\cD_H(\zeta)}\,
   \varepsilon_m.
 \label{eq:tpc61-defect-convex-combination}
\end{equation}
This is a convex combination of the row defects.  It is at most their
maximum, and concentrating \(\zeta\) on a maximizing row attains that
maximum after a harmless scalar rescaling into the coefficient class.
The represented identity follows from \(\rho_m=1-\varepsilon_m\), or by
the same argument with a minimum.
\end{proof}

For one distinguished vector, the exact and generally weaker quantity is
\begin{equation}
 \varepsilon_*(\zeta)
 :=\frac{\cD_M(\zeta)}{\cD_H(\zeta)}
 =\sum_m
   \frac{W_H(m)|\zeta_m|^2}{\cD_H(\zeta)}\varepsilon_m.
 \label{eq:tpc61-distinguished-defect}
\end{equation}
A favorable value of \(\varepsilon_*(\zeta)\) does not imply the uniform
bound in \eqref{eq:tpc61-sharp-missing-minimax}.

\begin{corollary}[Exact missing-only-row kill test]
\label{cor:tpc61-missing-only-row-kill}
If one active row satisfies
\begin{equation}
 W_R(m_0)=0<W_M(m_0),
 \label{eq:tpc61-missing-only-row}
\end{equation}
then the optimal uniform represented-mass constant is zero and the optimal
uniform missing fraction is one.
\end{corollary}

\subsection{What the inherited relation does not determine}

The papers preceding TPC--61 call \(\cI_X\) the retained or supported
pair relation, but impose no rowwise completion rule on it; arbitrary
finite cofactor multiplicity is explicitly allowed
\citep{WangTPC53,WangTPC54}.  Consequently the displayed inherited
hypotheses do not force
\begin{equation}
 \cI_m\cap\supp\nu_m\ne\varnothing
 \quad\text{on every active row}.
 \label{eq:tpc61-unproved-row-intersection}
\end{equation}

\begin{proposition}[Interface-insufficiency theorem]
\label{prop:tpc61-interface-insufficiency}
Within the finite-relation hypothesis class stated in
TPC--53--TPC--60, there is no universal constant \(c>0\), and no exponent
\(\sigma>0\), for which either
\begin{equation}
 \cD_R\ge c\cD_H
 \qquad\text{or}\qquad
 \cD_M\le X^{-\sigma+o(1)}\cD_H
 \label{eq:tpc61-unavailable-uniform-conclusions}
\end{equation}
follows from scale compatibility, finiteness of \(\cI_X\), and boundedness
of the row coefficients alone.
\end{proposition}

\begin{proof}
Those assumptions do not exclude a retained relation which omits every
realized cofactor on one otherwise active row.  Such a row satisfies
\eqref{eq:tpc61-missing-only-row}, and
\cref{cor:tpc61-missing-only-row-kill} gives both failures.  This is a
non-implication for the stated interface, not a claim that the upstream
physical provenance necessarily chooses such a relation.
\end{proof}

\begin{example}[Average-row dilution does not repair a uniform defect]
\label{ex:tpc61-average-row-trap}
Take any number of completed rows with \(\varepsilon_m=0\) and one active
missing-only row with \(\varepsilon_{m_0}=1\).  The unweighted average of
the row defects can tend to zero as the number of completed rows grows,
whereas \eqref{eq:tpc61-sharp-missing-minimax} remains equal to one.
The coefficient vector supported on \(m_0\) detects the bad row exactly.
Thus an average census can support a distinguished-vector statement only
after the physical coefficient weights in
\eqref{eq:tpc61-distinguished-defect} are retained.
\end{example}

The result above closes only a proof architecture: it says that the opaque
``retained relation'' interface is insufficient.  It is not a negative
theorem about the actual upstream carrier.  A genuine arithmetic advance
must either reopen the provenance of \(\cI_X\), or directly estimate the
realized weighted measures \(\nu_m(I_m^c)\).
